% ══════════════════════════════════════════════════════════════════════════════
% Lecture 15: Flows and Cuts – Max-Flow Min-Cut Theorem
% ══════════════════════════════════════════════════════════════════════════════

% ─────────────────────────────────────────────────────────────────────────────
\subsection{Full Example of Ford-Fulkerson}

To illustrate the method, here is the evolution of the flow on a network with 
integer capacities until the optimum $|f| = 29$:

\begin{center}
\begin{tikzpicture}[font=\small, >=Stealth,
    node/.style={circle, draw, fill=blue!8, minimum size=0.62cm, inner sep=1pt}]
  \node[node, fill=green!20] (s)  at (0,   1.5) {$s$};
  \node[node]                (a)  at (1.6,  3)  {$a$};
  \node[node]                (b)  at (1.6,  1.5){$b$};
  \node[node]                (c)  at (1.6,  0)  {$c$};
  \node[node]                (d)  at (3.2,  3)  {$d$};
  \node[node]                (e)  at (3.2,  1.5){$e$};
  \node[node]                (f)  at (3.2,  0)  {$f$};
  \node[node]                (g)  at (4.8,  3)  {$g$};
  \node[node]                (h)  at (4.8,  0)  {$h$};
  \node[node, fill=red!20]   (t)  at (6.4,  1.5){$t$};

  % Final optimal flow f/c
  \draw[->,thick] (s) -- node[above,font=\tiny]{10/10}(a);
  \draw[->,thick] (s) -- node[right,font=\tiny]{5/5} (b);
  \draw[->,thick] (s) -- node[below,font=\tiny]{14/15}(c);
  \draw[->,thick] (a) -- node[above,font=\tiny]{2/2} (d);
  \draw[->,thick] (a) -- node[above,font=\tiny]{9/9} (e);
  \draw[->,thick] (b) -- node[right,font=\tiny]{7/8} (e);
  \draw[->,thick] (c) -- node[below,font=\tiny]{4/4} (f);
  \draw[->,thick] (c) -- node[below,font=\tiny]{6/6} (e);
  \draw[->,thick] (d) -- node[above,font=\tiny]{3/15}(g);
  \draw[->,thick] (e) -- node[right,font=\tiny]{0/15}(g);
  \draw[->,thick] (e) -- node[right,font=\tiny]{0/15}(h);
  \draw[->,thick] (f) -- node[below,font=\tiny]{16/16}(h);
  \draw[->,thick] (g) -- node[above,font=\tiny]{9/10} (t);
  \draw[->,thick] (h) -- node[below,font=\tiny]{10/10}(t);
  \node[below=6pt, keyred, font=\small] at (3.2, -0.7)
        {Optimal flow: $|f| = 29$ — no augmenting path in $G_f$};
\end{tikzpicture}
\end{center}

% ══════════════════════════════════════════════════════════════════════════════
\section*{Cuts in Flow Networks}
\addcontentsline{toc}{subsection}{Cuts and Max-Flow Min-Cut Theorem}
% ══════════════════════════════════════════════════════════════════════════════

% ─────────────────────────────────────────────────────────────────────────────
\subsection{Definition of a Cut}

\begin{definition}{Cut}{cut}
A \textbf{cut} of the flow network $G = (V, E)$ is a partition of $V$ into
two sets $S$ and $T = V \setminus S$ such that $s \in S$ and $t \in T$.
\end{definition}

\begin{center}
\begin{tikzpicture}[font=\small, >=Stealth,
    node/.style={circle, draw, fill=blue!8, minimum size=0.65cm, inner sep=1pt}]
  \node[node, fill=green!20] (s)  at (0, 1)   {$s$};
  \node[node, fill=blue!20]  (v1) at (2, 2)   {$v_1$};
  \node[node, fill=blue!20]  (v2) at (2, 0)   {$v_2$};
  \node[node, fill=orange!20](v3) at (4, 2)   {$v_3$};
  \node[node, fill=orange!20](v4) at (4, 0)   {$v_4$};
  \node[node, fill=red!20]   (t)  at (6, 1)   {$t$};

  \draw[->, thick] (s)  -- node[above,font=\tiny]{11/16}(v1);
  \draw[->, thick] (s)  -- node[below,font=\tiny]{8/13} (v2);
  \draw[->, thick] (v1) -- node[above,font=\tiny]{12/12}(v3);
  \draw[->, thick] (v1) -- node[right,font=\tiny]{0/4}  (v2);
  \draw[->, thick] (v2) -- node[below,font=\tiny]{11/14}(v4);
  \draw[->, thick] (v2) -- node[left, font=\tiny]{0/9}  (v3);
  \draw[->, thick] (v3) -- node[above,font=\tiny]{19/20}(t);
  \draw[->, thick] (v4) -- node[below,font=\tiny]{7/7}  (t);
  \draw[->, thick] (v4) -- node[right,font=\tiny]{4/4}  (v3);

  % Cut boundary
  \draw[keyred, dashed, very thick] (3, -0.6) -- (3, 2.6)
        node[above, keyred]{cut $(S,T)$};
  \node[keyblue, font=\bfseries\tiny] at (1, 2.4) {$S = \{s, v_1, v_2\}$};
  \node[keyred,  font=\bfseries\tiny] at (5, 2.4) {$T = \{v_3, v_4, t\}$};
\end{tikzpicture}
\end{center}

% ─────────────────────────────────────────────────────────────────────────────
\subsection{Net Flow and Capacity of a Cut}

\begin{definition}{Net Flow and Capacity of a Cut}{cutflow}
The \textbf{net flow} through the cut $(S, T)$ is:
\[
  f(S, T) = \underbrace{\sum_{u \in S,\, v \in T} f(u,v)}_{\text{flow leaving }S}
            - \underbrace{\sum_{u \in S,\, v \in T} f(v,u)}_{\text{flow entering }S}.
\]

The \textbf{capacity} of the cut is:
\[
  c(S, T) = \sum_{u \in S,\, v \in T} c(u, v).
\]
\end{definition}

\textbf{Example} (cut $S = \{s, v_1, v_2\}$):
the net flow is $12 + 11 - 4 = 19$ and the capacity is $12 + 14 = 26$.

% ─────────────────────────────────────────────────────────────────────────────
\subsection{Net Flow is Always Equal to the Flow Value}

\begin{theorem}{Net Flow = Flow Value}{netflow}
For any cut $(S, T)$,
\[
  |f| = f(S, T).
\]
\end{theorem}

\begin{proof}
By induction on $|S|$.

\textbf{Base Case} ($S = \{s\}$): $f(\{s\}, T) = \sum_v f(s,v) - \sum_v f(v,s) = |f|$
by definition of the flow value.

\textbf{Inductive Step} ($S = S' \cup \{w\}$): when adding $w$ to $S'$,
the edges between $w$ and $S'$ move from the boundary to the interior of $S$.
The net flow variation is exactly
$\text{flow entering }w - \text{flow leaving }w = 0$
by flow conservation ($w \ne s, t$). Therefore $f(S, T) = f(S', T)$ and the value remains $|f|$.
\end{proof}

\textbf{Immediate Consequence:} since $f(u,v) \ge 0$ and $f(u,v) \le c(u,v)$,
\[
  |f| = f(S,T) \le \sum_{u \in S, v \in T} f(u,v) \le \sum_{u \in S, v \in T} c(u,v) = c(S,T).
\]

Thus \textbf{every flow is bounded by the capacity of any cut}, which implies:
\[
  \text{max-flow} \le \text{min-cut}.
\]

% ─────────────────────────────────────────────────────────────────────────────
\subsection{Max-Flow Min-Cut Theorem}

\begin{theorem}{Max-Flow Min-Cut}{maxflowmincut}
Let $G = (V, E)$ be a flow network with source $s$, sink $t$ and capacities $c$.
The following three assertions are equivalent:
\begin{enumerate}
  \item $f$ is a \textbf{maximum} flow.
  \item $G_f$ has \textbf{no augmenting path}.
  \item $|f| = c(S, T)$ for a \textbf{minimum} cut $(S, T)$.
\end{enumerate}
\end{theorem}

\begin{proof}
\textbf{$(1) \Rightarrow (2)$:} Suppose by contradiction that $G_f$ has
an augmenting path $p$. The Ford-Fulkerson method would augment $f$ along
$p$, increasing the flow value, which contradicts that $f$ is maximum.

\textbf{$(2) \Rightarrow (3)$:} Let $S = \{v \in V : \text{there exists a path from } s
\text{ to } v \text{ in } G_f\}$ and $T = V \setminus S$.
By assumption, $t \notin S$, so $(S,T)$ is a valid cut.

\begin{itemize}[noitemsep]
  \item \textbf{Every edge $(u,v)$ leaving $S$ in $G$} has $f(u,v) = c(u,v)$
        (saturated). Otherwise $c_f(u,v) = c(u,v) - f(u,v) > 0$ and $v$ would be reachable
        from $s$ in $G_f$, contradicting $v \notin S$.

  \item \textbf{Every edge $(v,u)$ entering $S$ from $T$ in $G$} has $f(v,u) = 0$.
        Otherwise $c_f(u,v) = f(v,u) > 0$ and $v$ would be reachable from $s$ in $G_f$.
\end{itemize}

\begin{center}
\begin{tikzpicture}[font=\small, >=Stealth,
    node/.style={circle, draw, fill=blue!8, minimum size=0.55cm, inner sep=1pt}]
  % Left: original graph with saturated/zero edges
  \begin{scope}[xshift=0cm]
    \node[draw, ellipse, fill=blue!10, minimum width=2cm, minimum height=1.4cm]
          at (0.4, 0) {};
    \node[font=\bfseries\small, keyblue] at (0.4, 0) {$S$};
    \node[draw, ellipse, fill=orange!10, minimum width=1.6cm, minimum height=1.4cm]
          at (3.0, 0) {};
    \node[font=\bfseries\small, keyred] at (3.0, 0) {$T$};
    \draw[->, keyred, very thick] (1.4, 0.3) -- node[above,font=\tiny,keyred]{$f=c$ (saturated)} (2.2, 0.3);
    \draw[->, keygreen, very thick] (2.2, -0.3) -- node[below,font=\tiny,keygreen]{$f=0$} (1.4, -0.3);
    \node[below=14pt] at (1.8, -0.6) {\textit{graph $G$}};
  \end{scope}
  % Right: residual network
  \begin{scope}[xshift=5.5cm]
    \node[draw, ellipse, fill=blue!10, minimum width=2cm, minimum height=1.4cm]
          at (0.4, 0) {};
    \node[font=\bfseries\small, keyblue] at (0.4, 0) {$S$};
    \node[draw, ellipse, fill=orange!10, minimum width=1.6cm, minimum height=1.4cm]
          at (3.0, 0) {};
    \node[font=\bfseries\small, keyred] at (3.0, 0) {$T$};
    \draw[->, keyred, very thick] (2.2, 0.3) -- node[above,font=\tiny,keyred]{$c_f = c-c = 0$} (1.4, 0.3);
    \node[below=14pt] at (1.8, -0.6) {\textit{residual network $G_f$}};
    \node[gray, font=\tiny] at (1.8, -0.1) {(no edge toward $T$)};
  \end{scope}
\end{tikzpicture}
\end{center}

Therefore:
\[
  |f| = f(S,T) = \sum_{\substack{u\in S\\v\in T}} f(u,v) - \sum_{\substack{u\in S\\v\in T}} f(v,u)
       = \sum_{\substack{u\in S\\v\in T}} c(u,v) - 0 = c(S,T).
\]

\textbf{$(3) \Rightarrow (1)$:} Since $|f| \le c(S,T)$ for any cut,
if $|f| = c(S,T)$ for a minimum cut, the flow cannot be improved. Thus $f$ is maximum.
\end{proof}

% ─────────────────────────────────────────────────────────────────────────────
\subsection{Illustration on the Example}

In the 7-node network (optimal flow $|f| = 29$), the minimum cut identified by
$S = \{s\} \cup \{\text{nodes reachable from }s\text{ in }G_f\}$ satisfies
$c(S,T) = 29$: the edges of the cut are exactly saturated, and no augmenting
path exists in $G_f$.

% ─────────────────────────────────────────────────────────────────────────────
\subsection{Lecture Summary}

\begin{tcolorbox}[colback=lightblue, colframe=keyblue, fonttitle=\bfseries,
                  title={Key Points – Flows and Cuts}]
\begin{itemize}[noitemsep]
  \item \textbf{Cut $(S,T)$}: partition $V = S \cup T$ with $s \in S$, $t \in T$.
  \item \textbf{Net Flow} $f(S,T) = \sum_{S \to T} f - \sum_{T \to S} f$.
  \item \textbf{Capacity} $c(S,T) = \sum_{u\in S, v\in T} c(u,v)$.
  \item \textbf{Fundamental Identity:} $|f| = f(S,T)$ for any cut,
        hence $|f| \le c(S,T)$.
  \item \textbf{Max-Flow Min-Cut Theorem:}
    \[
      \max |f| = \min c(S,T).
    \]
    Three equivalent conditions: max flow $\Leftrightarrow$ no augmenting
    path in $G_f$ $\Leftrightarrow$ flow equals min-cut capacity.
  \item \textbf{Ford-Fulkerson returns an optimal flow} because it stops
        precisely when condition (2) is verified.
\end{itemize}
\end{tcolorbox}